\begin{abstract}
Large-scale text-to-image diffusion models have demonstrated remarkable abilities in producing photorealistic visuals from textual descriptions. Effectively steering these powerful models for diverse downstream applications, however, remains a key open question. In response, we propose Orthogonal Finetuning~(OFT), a systematic finetuning strategy designed to adapt text-to-image diffusion models for specific downstream scenarios. Distinct from prevailing approaches, OFT can rigorously maintain the hyperspherical energy, which reflects the relational structure among neurons constrained to the unit hypersphere. Our analysis indicates that conserving this energy is fundamental for sustaining the models' semantic generation capabilities. To enhance stability during finetuning, we further introduce Constrained Orthogonal Finetuning (COFT), which applies an explicit constraint to control the radius of the hypersphere. In particular, we examine two representative finetuning challenges: subject-driven generation, aiming to synthesize novel images of a specific entity given a handful of reference images and a text prompt, and controllable generation, where the objective is to incorporate auxiliary control signals. Our empirical evaluations demonstrate that the OFT framework surpasses current techniques both in terms of generation fidelity and training efficiency.
\end{abstract}

\section{Introduction}

The latest advancements in text-to-image diffusion models~\cite{saharia2022photorealistic,ramesh2022hierarchical,rombach2022high} have established new benchmarks for text-driven, high-quality image creation. Nonetheless, even with notable performance, textual prompts alone frequently lack the expressiveness to achieve detailed and precise manipulations in generated outputs. To overcome these limitations, this work addresses two principal categories of text-to-image generation:

\begin{itemize}[leftmargin=*,nosep]

    \item \textbf{Subject-driven generation}~\cite{ruiz2023dreambooth}: With a limited set of images depicting a target subject, the objective is to generate new imagery portraying this subject in varied contexts, guided by textual prompts.
    \item \textbf{Controllable generation}~\cite{zhang2023adding,mou2023t2i}: By supplying supplementary control cues (\emph{e.g.}, canny edge maps, segmentation masks), the goal is to condition image synthesis on both these signals and accompanying textual instructions.
\end{itemize}

At their core, both challenges revolve around developing strategies for fine-tuning text-to-image diffusion models in such a way that the generative capabilities acquired during pretraining are not compromised. The requirements for an optimal finetuning approach can be outlined as follows: (1) \emph{training efficiency}, characterized by minimizing both the number of parameters to be updated and the required training epochs; and (2) \emph{generalizability preservation}, which refers to maintaining the model’s capacity for high-quality and diverse generation. Approaches to finetuning usually fall into two categories: directly updating neuron weights with a modest learning rate (\emph{e.g.}, \cite{ruiz2023dreambooth}), or incorporating an auxiliary module consisting of re-parameterized weights (\emph{e.g.}, \cite{hulora2022,zhang2023adding}). Although these methods are straightforward, neither can reliably ensure the retention of generative abilities established during pretraining. Additionally, there remains a conceptual gap in the literature regarding how to systematically construct effective finetuning techniques or to determine appropriate values for critical hyperparameters, such as training duration. Central to these challenges is the absence of a rigorous metric for evaluating to what extent the finetuned model preserves the generative competence of its pretrained counterpart. Prevailing methods tend to presume that minimizing the Euclidean distance in parameter space between the finetuned and pretrained models equates to better preservation of generative performance. Consequently, finetuning protocols often employ extremely small learning rates. However, while this perspective is sometimes valid, we contend that Euclidean proximity alone cannot adequately capture the nuances of semantic preservation. As such, adopting a more structurally informed measure to assess the relationship between the finetuned and pretrained models could substantially enhance both the fidelity of pretraining retention and the overall robustness of the finetuning process.

Building on empirical findings that hyperspherical similarity effectively encodes semantic relationships~\cite{liu2017deep,liu2018decoupled,chen2020angular}, we utilize the concept of hyperspherical energy~\cite{liu2018learning} to represent the relational structure among neurons within a layer. Hyperspherical energy is calculated as the aggregated hyperspherical similarity (\emph{e.g.}, cosine similarity) across all neuron pairs in the same layer, thus quantifying the uniform distribution of neurons on a unit hypersphere~\cite{liu2021learning}. Our working hypothesis posits that an optimally finetuned model should exhibit minimal alteration in hyperspherical energy when compared with its pretrained counterpart. A straightforward strategy would involve applying a regularization term to enforce invariance of hyperspherical energy throughout the finetuning process, but this approach does not ensure that the discrepancy in hyperspherical energy is adequately reduced. To address this, we exploit a key invariance property: pairwise hyperspherical similarity remains unchanged under a shared orthogonal transformation applied to all neurons. This motivates the design of Orthogonal Finetuning~(OFT), a method that adapts large text-to-image diffusion models for specific tasks while maintaining unchanged hyperspherical energy. The core mechanism of OFT involves learning a shared orthogonal transformation per layer that preserves pairwise angular relationships among neurons. Alternatively, OFT can be interpreted as redefining the coordinate basis for neurons in a given layer. Taking into account the observation that smaller Euclidean distances between the parameters of finetuned and pretrained models correlate with better preservation of pretrained capabilities, we further introduce Constrained Orthogonal Finetuning~(COFT), a variant of OFT which restricts the finetuned parameters to remain within a hypersphere of fixed radius around the pretrained configuration.

The rationale behind employing orthogonal transformations in neuron finetuning is partly inspired by concepts from the 2D Fourier transform, which allows an image to be represented through its magnitude and phase spectra. The phase spectrum, reflecting the angular relationship between the input and its basis, retains the core semantic content. Notably, reconstructing an image using its phase spectrum and a random magnitude spectrum largely preserves the original semantic features. This observation indicates that adjusting the orientation of neurons is essential for implementing semantic changes in generated images—central to both subject-driven and controllable generation tasks. Nevertheless, if neuron orientations are altered too freely, the integrity of the model’s pretrained generative capabilities is likely to be compromised. To address this, we propose to maintain the relative angles between all neuron pairs, based on the premise that these angles are integral to how neural networks encode knowledge. This perspective motivates our approach: we introduce layer-wise orthogonal transformations, shared across neurons, so as to ensure that the hyperspherical energy is invariant throughout the process.

Our approach also draws from the methodology of orthogonal over-parameterized training~\cite{liu2021orthogonal}, where neural networks for classification are trained by applying orthogonal transformations to an initially random network. The theoretical justification here is that random initialization typically yields low hyperspherical energy on average; \cite{liu2021orthogonal} demonstrates that preserving low hyperspherical energy during training (which has been shown to correlate with improved generalization~\cite{liu2018learning,lin2020regularizing}) supports the formation of robust classification models. While \cite{liu2021orthogonal} validates the capacity of orthogonal transformations to yield well-generalizing classifiers, our focus shifts to finetuning text-to-image diffusion models with the goals of enhanced controllability and superior downstream generative outcomes. Key distinctions exist between OFT and \cite{liu2021orthogonal}. First, whereas \cite{liu2021orthogonal} seeks to further reduce hyperspherical energy, OFT is designed to maintain the energy level established by the pretrained model, thereby safeguarding the intrinsic learned structure during finetuning. Attempts to minimize hyperspherical energy in diffusion model finetuning could, in fact, damage the pre-existing semantic representations. Second, OFT incorporates an explicit constraint to minimize deviations from the pretrained parameters, resulting in a constrained form, whereas \cite{liu2021orthogonal} imposes no such restrictions. Striking an optimal balance between adaptability and preservation of the original model is central to effective finetuning, and we contend that the OFT framework succeeds in meeting this requirement. Our main contributions are summarized as follows:

\begin{itemize}[leftmargin=*,nosep]

    \item We introduce a new finetuning approach—Orthogonal Finetuning—that steers text-to-image diffusion models towards enhanced controllability. Additionally, we develop a constrained version that restricts the angular shift relative to the original pretrained model to further bolster finetuning stability.
    \item Unlike previous finetuning strategies, OFT adapts the model while rigorously maintaining the hyperspherical energy. Our empirical results suggest that preserving this property is vital for retaining the semantic generative capacities of the original model.
    \item We utilize OFT in both subject-driven and controllable generation scenarios. Through extensive experiments, we show that OFT considerably outperforms existing techniques in generation fidelity, speed of convergence, and stability during finetuning. Notably, OFT is also more sample-efficient, requiring far fewer training examples and epochs to reach strong performance.
    \item For the task of controllable generation, we present a novel control consistency metric to systematically assess controllability. This metric operates by inferring the control signal from an output image and comparing it against the original input control, providing additional evidence for OFT’s robust controllability.
\end{itemize}

\section{Related Work}

\textbf{Text-to-image diffusion models}. The landscape of text-to-image synthesis has rapidly advanced~\cite{saharia2022photorealistic,ramesh2022hierarchical,rombach2022high,gu2022vector,nichol2022glide}, thanks to strides in diffusion-based generative modeling~\cite{ho2020denoising,song2021score,song2021denoising,dhariwal2021diffusion} and the evolution of vision-language learning methods~\cite{su2020vlbert,lu2019vilbert,radford2021learning,wang2021simvlm,li2022blip,li2023blip,alayrac2022flamingo,schuhmann2022laion}. Approaches such as GLIDE~\cite{nichol2022glide} and Imagen~\cite{saharia2022photorealistic} apply diffusion in pixel space; GLIDE aligns the text encoder and diffusion process jointly via paired text-image training, while Imagen employs a frozen, pretrained text encoder. Latent space diffusion is leveraged by models like Stable Diffusion~\cite{rombach2022high} and DALL-E2~\cite{ramesh2022hierarchical}, with Stable Diffusion relying on VQ-GAN~\cite{esser2021taming} for latent visual codes, and DALL-E2 using CLIP~\cite{radford2021learning} to build a shared embedding space for both modalities. Beyond diffusion, other generative frameworks, such as generative adversarial networks~\cite{reed2016generative,zhang2017stackgan,xu2018attngan,li2019controllable} and autoregressive techniques~\cite{ramesh2021zero,ding2021cogview,wu2022nuwa,yuscaling}, have been applied to this domain. OFT is designed to be broadly compatible and can be integrated with any text-to-image diffusion model.

\textbf{Subject-driven generation}. To address the issue of subject alteration, prior works~\cite{avrahami2022blended,nichol2022glide} incorporate a user-provided mask as supplementary input. Alternative methods, such as inversion~\cite{choi2021ilvr,dhariwal2021diffusion,ramesh2022hierarchical,gal2022image}, adjust context while retaining the subject. Techniques in~\cite{hertz2022prompt} facilitate both localized and holistic edits, even without mask inputs. However, these approaches often fall short of faithfully preserving subject identity details. Pivotal Tuning~\cite{roich2022pivotal} circumvents this by locally updating a generator near an inverted latent code, combined with identity-preserving regularization. Similarly,~\cite{nitzan2022mystyle} tailors a generative face prior to a specific individual's image set, and~\cite{casanova2021instance} enables instance-level variation though occasionally neglects distinctive instance traits. DreamBooth~\cite{ruiz2023dreambooth} customizes a diffusion model to a new subject using a dedicated token and a limited number of images, with both reconstruction and class-specific prior losses for guidance. While OFT utilizes the DreamBooth paradigm, it distinguishes itself by implementing orthogonal updates during finetuning rather than standard parameter adjustments with reduced learning rates.

\textbf{Controllable generation}. Image-to-image translation can be interpreted as a specialized case of controllable image generation. Historically, conditional generative adversarial networks have been predominantly utilized for this purpose~\cite{isola2017image,zhu2017toward,wang2018high,park2019semantic,choi2018stargan}. Recent advancements also demonstrate the application of diffusion models in image-to-image translation tasks~\cite{saharia2022palette,voynov2022sketch,wang2022pretraining}. Notably, ControlNet~\cite{zhang2023adding} introduces a method to steer a pretrained diffusion model by finetuning it in response to supplementary control inputs, achieving significant improvements in controllable image synthesis. Similarly, the contemporaneous T2I-Adapter~\cite{mou2023t2i} enhances image controllability by adapting a pretrained diffusion model via finetuning. Consistent with the task formulations outlined in \cite{zhang2023adding,mou2023t2i}, we utilize OFT to refine pretrained diffusion models, which results in improved controllability, requiring less annotated data and fewer additional finetuning parameters. Importantly, OFT maintains computational efficiency at inference time and does not impose extra overhead.

\textbf{Model finetuning}. Recently, finetuning large-scale pretrained models for downstream applications has seen a surge in adoption~\cite{devlin2018bert,brown2020language,he2022masked}. Adaptation-based strategies, such as those in \cite{hulora2022,houlsby2019parameter,pfeiffer2020adapterfusion}, have been extensively explored in the field of natural language processing as a subclass of finetuning methods. The LoRA approach~\cite{hulora2022}, which is particularly related to OFT, posits a low-rank approximation for incremental weight modifications during finetuning. Unlike LoRA, OFT leverages a shared orthogonal transformation across layers to multiplicatively modify neuron weights. This approach has theoretical guarantees for preserving the pairwise angles between neurons within a layer, thereby enhancing model stability.

\section{Orthogonal Finetuning}

\subsection{Why Does Orthogonal Transformation Make Sense?}

To clarify the motivation for using orthogonal transformations during the finetuning of text-to-image diffusion models, we break down the rationale into two core aspects: (1) the benefits of adjusting neuron angles (i.e., their direction), and (2) the justification for using orthogonal transformations as the tool for modifying these angles.

Addressing the first question, our approach is informed by the findings in \cite{liu2018decoupled,chen2020angular}, which suggest that differences in angular features are highly indicative of the semantic gap. SphereNet~\cite{liu2017deep} demonstrates that constraining all neurons to a unit hypersphere during neural network training results in similar expressive power and often superior generalization ability, suggesting that the directional components of neurons are critical for capturing salient data features. To highlight the significance of neuron directions, we carry out a simple experiment, depicted in Figure~\ref{fig:toy_ae}, where a standard convolutional autoencoder is trained on flower image data. During training, feature maps are computed via the standard inner product ($z$ representing the output element from convolution with kernel $\bm{w}$ applied to input $\bm{x}$ in a sliding window). In the evaluation phase, we compare three different schemes for feature map computation: (a) the inner product utilized during training, (b) retaining only magnitude information, and (c) using only angular information. Figure~\ref{fig:toy_ae} illustrates that neuron angular components are nearly sufficient to reconstruct the original images, whereas neuron magnitudes do not carry meaningful information. It is important to note that the angular activation (c) is applied only at inference time; the network was trained exclusively with the inner product. This outcome suggests that the angles, or directions, of neurons encode most of the semantic content in the input data. Therefore, adjusting the directions of neurons offers an effective means for modifying the semantic attributes of images.

With respect to the second question, the most straightforward method for adjusting neuron directions is to jointly rotate or reflect all neurons in a given layer, naturally invoking an orthogonal transformation. While other angular modifications—such as independently rotating neurons by different angles—could offer greater flexibility, we find orthogonal transformations provide a valuable trade-off between adaptability and structure. Furthermore, \cite{liu2021orthogonal} provides evidence that orthogonal transformations are sufficiently expressive for training neural networks. To bolster our position, we conduct a controlled experiment following the setup of ControlNet~\cite{zhang2023adding} (results shown in Figure~\ref{fig:ortho}), to examine the impact of enforcing orthogonality as a regularization mechanism. The results demonstrate that our standard OFT achieves stable performance and precise control upon convergence (epoch 20). By contrast, removing the orthogonality constraint results in an inability to produce realistic images or achieve intended control. These findings affirm the crucial role of the orthogonality constraint in OFT.

\subsection{General Framework}

The standard approach to finetuning involves applying gradient descent with a modest learning rate to update the model parameters (or selected layers thereof). Utilizing a small learning rate acts as a regularizer, effectively limiting the departure from the pretrained weights; thus, conventional finetuning can be viewed as seeking to minimize $\thickmuskip=2mu \medmuskip=2mu \| \bm{M} - \bm{M}^0\|$, where $\bm{M}$ denotes the weights after finetuning and $\bm{M}^0$ those before. Although this soft restriction is conceptually reasonable, it may not sufficiently constrain parameter updates when adapting large-scale models. To mitigate this, LoRA incorporates an explicit low-rank condition on the parameter shift, imposing $\thickmuskip=2mu \medmuskip=2mu \text{rank}(\bm{M} - \bm{M}^0)=r'$ for a small, user-chosen $r'$. Unlike LoRA, OFT enforces constraints at the level of neuron interactions by requiring $\thickmuskip=2mu \medmuskip=2mu \| \text{HE}(\bm{M})-\text{HE}(\bm{M}^0) \|=0$, with $\text{HE}(\cdot)$ representing the hyperspherical energy corresponding to a given weight matrix.

To illustrate, consider a fully connected layer described by the weight matrix $\thickmuskip=2mu \medmuskip=2mu \bm{W}=\{\bm{w}_1,\cdots,\bm{w}_n\}\in\mathbb{R}^{d\times n}$, where each $\bm{w}_i\in\mathbb{R}^{d}$ constitutes an individual neuron and $\bm{W}^0$ contains the pretrained weights. Given an input vector $\thickmuskip=2mu \medmuskip=2mu \bm{x}\in\mathbb{R}^d$, the corresponding output $\thickmuskip=2mu \medmuskip=2mu \bm{z}\in\mathbb{R}^n$ is computed as $\thickmuskip=2mu \medmuskip=2mu \bm{z}=\bm{W}^\top\bm{x}$. The OFT framework can thus be cast as the following optimization, which encourages minimal change in hyperspherical energy relative to the original parameters:

\begin{equation}\label{eq:oft_principle}
\footnotesize
\min_{\bm{W}} \left\| \text{HE}(\bm{W})-\text{HE}(\bm{W}^0)\right\|~~\Leftrightarrow~~\min_{\bm{W}} \bigg{\|} \sum_{i\neq j} \|\hat{\bm{w}}_i-\hat{\bm{w}}_j\|^{-1}-\sum_{i\neq j} \|\hat{\bm{w}}^0_i-\hat{\bm{w}}^0_j\|^{-1}\bigg{\|}
\end{equation}

Here, each normalized neuron is represented as $\thickmuskip=2mu \medmuskip=2mu \hat{\bm{w}}_i:=\bm{w}_i/\|\bm{w}_i\|$, and the hyperspherical energy of $\bm{W}$ is given by $\thickmuskip=2mu \medmuskip=2mu \text{HE}(\bm{W}):= \sum_{i\neq j} \|\hat{\bm{w}}_i-\hat{\bm{w}}_j\|^{-1}$. It is straightforward to see that the optimum for Eq.~\eqref{eq:oft_principle} is zero, and this is attainable when $\bm{W}$ and $\bm{W}^0$ differ only by an orthogonal transformation; formally, if $\thickmuskip=2mu \medmuskip=2mu \bm{W}=\bm{R}\bm{W}^0$ with $\thickmuskip=2mu \medmuskip=2mu \bm{R}\in\mathbb{R}^{d\times d}$ an orthogonal matrix (where determinant $1$ or $-1$ indicates a rotation or a reflection, respectively), then the constraint is met. The central notion of OFT, therefore, is to restrict finetuning to transformations parameterized by layer-wise orthogonal matrices that alter neuron representations within each layer.

\begin{equation}\label{eq:oft_general}
    \footnotesize
    \bm{z}=\bm{W}^\top\bm{x}=(\bm{R}\cdot\bm{W}^0)^\top\bm{x},~~~\text{s.t.}~\bm{R}^\top\bm{R}=\bm{R}\bm{R}^\top=\bm{I}
\end{equation}

In this equation, $\bm{W}$ represents the weight matrix obtained after applying OFT fine-tuning, and $\bm{I}$ indicates the identity matrix. Figure~\ref{fig:block} provides an overview of OFT. To mirror the effect of zero initialization in LoRA, it is necessary for OFT to ensure that fine-tuning commences precisely from the original pretrained weights $\bm{W}^0$. This is accomplished by initializing the orthogonal matrix $\bm{R}$ as an identity matrix, guaranteeing that the finetuned parameters initially match the pretrained configuration. To maintain the orthogonality of $\bm{R}$ during optimization, one may employ differentiable orthogonalization techniques, such as those outlined in \cite{liu2021orthogonal,lezcano2019cheap}. In the following, we explore computationally tractable ways to enforce this orthogonality constraint.

\subsection{Efficient Orthogonal Parameterization}\label{sect:eff_ortho}

Although differentiable orthogonalization procedures like the Gram-Schmidt process can be utilized, they typically impose a prohibitive computational burden in large-scale scenarios~\cite{liu2021orthogonal}. For practical efficiency, we instead utilize Cayley parameterization, which facilitates the construction of orthogonal matrices. Concretely, we define the orthogonal matrix as $\thickmuskip=2mu \medmuskip=2mu \bm{R}=(\bm{I}+\bm{Q})(I-\bm{Q})^{-1}$, where $\bm{Q}$ is required to be skew-symmetric, i.e., $\thickmuskip=2mu \medmuskip=2mu \bm{Q}=-\bm{Q}^\top$. This strategy trades off generality for computational speed, as Cayley parameterization generates only orthogonal matrices with a determinant of $1$, meaning it covers only the special orthogonal group. Nevertheless, our experiments indicate this restriction does not adversely impact empirical results. However, even with Cayley parameterization, when the dimensionality $d$ is large, directly parameterizing $\bm{R}$ remains parameter inefficient. To mitigate this, we propose modeling $\bm{R}$ as a block-diagonal matrix partitioned into $r$ blocks, resulting in the structure:

\begin{equation}
    \footnotesize
    \bm{R}=\text{diag}(\bm{R}_1,\bm{R}_2,\cdots,\bm{R}_r)=\begin{bmatrix}
    \bm{R}_1\in \text{O}(\frac{d}{r}) &  &\\
    &\ddots & \\
    & & \bm{R}_r\in \text{O}(\frac{d}{r})
    \end{bmatrix}\in \text{O}(d)
\end{equation}

Here, $\text{O}(d)$ refers to the set of $d$-dimensional orthogonal matrices, where $\thickmuskip=2mu \medmuskip=2mu \bm{R}\in\mathbb{R}^{d\times d}$ as well as each block $\thickmuskip=2mu \medmuskip=2mu \bm{R}_i\in\mathbb{R}^{d/r\times d/r},~\forall i$, are orthogonal by definition. If we set $\thickmuskip=2mu \medmuskip=2mu r=1$, the block-diagonal structure reduces to a typical unconstrained orthogonal matrix. An orthogonal matrix of dimension $\thickmuskip=2mu \medmuskip=2mu d\times d$ is parameterized by $\thickmuskip=2mu \medmuskip=2mu d(d-1)/2$ degrees of freedom, leading to computational and parameter costs of order $\mathcal{O}(d^2)$. For an orthogonal matrix in block-diagonal form with $r$ blocks, the total number of free parameters is $\thickmuskip=2mu \medmuskip=2mu d(d/r-1)/2$, which reduces complexity to $\mathcal{O}(d^2/r)$. Further parameter sharing between the blocks, that is, enforcing $\thickmuskip=2mu \medmuskip=2mu\bm{R}_i = \bm{R}_j~\forall i\neq j$, brings the count down to $\mathcal{O}(d^2/r^2)$. Crucially, these reductions preserve the orthogonality of $\bm{R}$, so the ability to conserve hyperspherical energy is unaffected.

Let us compare the parameter usage of OFT with that of LoRA. LoRA with a low-rank parameter $r'$ contains $\thickmuskip=2mu \medmuskip=2mu r'(d+n)$ learnable parameters. Assuming both $r$ and $r'$ scale with the model's hidden size $d$ (e.g., $\thickmuskip=2mu \medmuskip=2mu r = r' = \alpha d$ with $\thickmuskip=2mu \medmuskip=2mu 0<\alpha\leq 1$), the parameter count for LoRA is $\mathcal{O}(d^2 + dn)$, while that for OFT is merely $\mathcal{O}(d)$. To provide a specific illustration: if the weight matrix is of size $\thickmuskip=2mu \medmuskip=2mu 128\times 128$, then with $\thickmuskip=2mu \medmuskip=2mu r' = 8$, LoRA contains $2{,}048$ trainable weights, while OFT, with $\thickmuskip=2mu \medmuskip=2mu r = 8$ (and no block sharing), involves only $960$.

\subsection{Constrained Orthogonal Finetuning}

It is possible to tighten the adaptation space for the finetuned model under OFT by confining it to a limited vicinity around the original parameters. The COFT method enforces this idea, performing inference according to

\begin{equation}\label{eq:coft_general}
    \footnotesize
    \bm{z} = \bm{W}^\top \bm{x} = (\bm{R}\cdot\bm{W}^0)^\top\bm{x},~~~\text{s.t.}~\bm{R}^\top \bm{R} = \bm{R}\bm{R}^\top = \bm{I},~\norm{\bm{R}-\bm{I}}\leq \epsilon
\end{equation}

where the matrix $\bm{R}$ is forced both to be orthogonal and to be at most $\epsilon$-distance from the identity. While the orthogonality requirement is readily satisfied using the Cayley parameterization (see Section~\ref{sect:eff_ortho}), enforcing the additional $\epsilon$-closeness constraint poses nontrivial challenges. To build intuition for the Cayley transformation, one can approximate $\thickmuskip=2mu \medmuskip=2mu \bm{R}=(\bm{I}+\bm{Q})(I-\bm{Q})^{-1}$ via its Neumann series expansion, yielding $\thickmuskip=2mu \med

series converges in the operator norm). Thus, the restriction $\thickmuskip=2mu \medmuskip=2mu\|\bm{R}-\bm{I}\|\leq\epsilon$ can be embedded within the Cayley transform. This yields an alternative constraint: $\thickmuskip=2mu \medmuskip=2mu \|\bm{Q}-\bm{0}\|\leq\epsilon'$, where $\bm{0}$ stands for the zero matrix, and $\epsilon'$ is a distinct hyperparameter governing error tolerance (not the same as $\epsilon$). Enforcing this new condition on the matrix $\bm{Q}$ can be done conveniently via projected gradient descent. Identity initialization of the orthogonal matrix $\bm{R}$ is attained by initializing $\bm{Q}$ as the zero matrix. COFT, therefore, can be regarded as an intersection of two explicit criteria: minimization of hyperspherical energy differences and a restriction on divergence from the pretrained parameters. While the latter is usually an implicit feature of standard finetuning techniques, COFT formalizes it as a specific constraint. Although OFT achieves strong results, we observe that COFT achieves even greater stability during finetuning, an improvement attributable to this additional deviation constraint. Figure~\ref{fig:eps_coft} illustrates the influence of $\epsilon$ on COFT's outcomes. The figure demonstrates that $\epsilon$ modulates the adaptability during finetuning: larger values of $\epsilon$ lead the COFT-adjusted model to behave more like one finetuned using OFT, whereas decreasing $\epsilon$ causes the COFT model to align more closely with the initial pretrained text-to-image diffusion model.

\subsection{Re-scaled Orthogonal Finetuning}

We introduce an uncomplicated extension of the conventional OFT by including a learnable scaling parameter for each neuron. The justification stems from the observation that scaling neuron outputs does not affect hyperspherical energy, since normalization negates the scale factor. Concretely, the forward computation is: $\thickmuskip=2mu \medmuskip=2mu\bm{z}=(\bm{R}\bm{W}^0\bm{D})^\top\bm{x}$\footnote{\textbf{Errata}: The NeurIPS camera-ready manuscript inaccurately stated the re-scaled OFT forward computation as $\thickmuskip=2mu \medmuskip=2mu\bm{z}=(\bm{D}\bm{R}\bm{W}^0)^\top\bm{x}$. The code implementation, however, was correct, leaving the findings in Appendix~\ref{app:rescaled} unaffected.}, where $\thickmuskip=2mu \medmuskip=2mu \bm{D}=\text{diag}(s_1,\cdots,s_n)\in\mathbb{R}^{n\times n}$ is a diagonal matrix containing trainable entries $s_1,\ldots,s_n$ with all $s_i>0$. Unlike the original OFT approach in Eq.~\eqref{eq:oft_general}, which involves learning only $\bm{R}$, both $\bm{D}$ and $\bm{R}$ are optimized here. This re-scaled OFT formulation enhances OFT’s expressiveness, introducing only a marginal number of extra learnable parameters. We adopt the baseline OFT setup in our experiments to emphasize the impact of orthogonal transformation alone, yet our findings suggest that the re-scaled variant generally provides additional benefits (refer to Appendix~\ref{app:rescaled}).

\section{Intriguing Insights and Discussions}

\textbf{OFT's architectural generality.} OFT is theoretically applicable to any neural architecture. For Transformer-based models, LoRA is customarily deployed on attention layer weights~\cite{hulora2022}. To ensure a fair comparison with LoRA, we restrict OFT in our experiments to finetuning attention weights exclusively. Beyond dense layers, OFT is particularly advantageous for adjusting convolutional layers, as the block-diagonal configuration of $\bm{R}$ gives rise to meaningful interpretations in this setting—unlike LoRA. For instance, with the number of blocks set equal to the count of input channels, each block is responsible for transforming a single neuron channel, echoing the concept of depth-wise convolutions~\cite{chollet2017xception}. In the case where all blocks of $\bm{R}$ are identical, the orthogonal matrix acts uniformly across all channels. An initial investigation into convolutional layer finetuning with OFT is provided in Appendix~\ref{app:convolution}.

\textbf{Connection to LoRA}. LoRA utilizes a low-rank matrix addition to limit substantial deviations in the pretrained weights. In comparison, OFT enforces orthogonality (full-rank) in the transformation applied to the pretrained weight matrix, thereby maintaining the integrity of the pretrained information. The forward computation in OFT can be expressed as $\thickmuskip=2mu \medmuskip=2mu \bm{z}=(\bm{R}\bm{W}^0)^\top\bm{x}=(\bm{W}^0+(\bm{R}-\bm{I})\bm{W}^0)^\top\bm{x}$, where $\thickmuskip=2mu \medmuskip=2mu (\bm{R}-\bm{I})\bm{W}^0$ serves a role analogous to the low-rank adjustment in LoRA. Since $\bm{W}^0$ is generally full-rank, OFT similarly facilitates a low-rank weight update if $\thickmuskip=2mu \medmuskip=2mu \bm{R}-\bm{I}$ is itself low-rank. LoRA introduces a tunable rank parameter $r'$, while OFT employs a block-diagonal size parameter $r$ to manage the count of learnable parameters. Notably, LoRA and OFT offer fundamentally different approaches for parameter-efficient tuning: LoRA leverages low-rank decomposition to reduce trainable parameter count, while OFT instead exploits sparsity through block-diagonal orthogonal structures.

\textbf{Why OFT converges faster}. As illustrated in Figure~\ref{fig:toy_ae}, updating neuron directions is the most direct means of altering their semantics, which is precisely the operation OFT targets. Furthermore, OFT can be interpreted as tuning parameters along a smooth hypersphere, resulting in a more favorable optimization landscape. This property is further corroborated by empirical evidence in \cite{liu2021orthogonal}.

\textbf{Why not minimize hyperspherical energy}. Unlike the approach in \cite{liu2021orthogonal}, our method does not focus on minimizing hyperspherical energy. In the context of classification, avoiding redundancy among neurons is important. Minimizing hyperspherical energy spaces neurons evenly on the hypersphere, but this objective is not suitable for finetuning since it risks undermining the valuable information captured during pretraining.

\textbf{Balancing flexibility and regularity in finetuning}. Our analysis reveals an inherent tension between model flexibility and imposed regularity. Conventional finetuning offers maximal adaptability but often lacks stability, frequently leading to phenomena such as model collapse. In contrast, OFT, despite its straightforward nature, achieves a well-calibrated compromise by maintaining pairwise angular relationships between neurons. Additionally, the block-diagonal structure in parameterization can be interpreted as a form of stronger regularization imposed on the orthogonal matrix.

\textbf{No increase in inference latency}. Distinct from ControlNet, the OFT framework imposes zero additional inference-time costs on the finetuned network. At inference, one can simply integrate the trained orthogonal matrix $\bm{R}$ with the pretrained weights $\bm{W}^0$, resulting in a combined weight matrix $\thickmuskip=2mu \medmuskip=2mu \bm{W}=\bm{R}\bm{W}^0$. Consequently, the model maintains the same inference speed as the original pretrained network.

\section{Experiments and Results}

\textbf{General experimental protocol}. All experiments are conducted using Stable Diffusion v1.5~\cite{rombach2022high} as the underlying text-to-image backbone. To maintain experimental fairness, sample images are randomly drawn from each method under comparison. For tasks involving subject-driven synthesis, our approach aligns with procedures in DreamBooth~\cite{ruiz2023dreambooth}; for controllable image generation, protocols adhere to those outlined in ControlNet~\cite{zhang2023adding} and T2I-Adapter~\cite{mou2023t2i}. To ensure parity with LoRA, both OFT and COFT are applied to identical model layers as LoRA. Additional experimental results and methodological specifics are documented in Appendix~\ref{app:settings}.

\subsection{Subject-driven Generation}

\textbf{Experimental setup}. DreamBooth~\cite{ruiz2023dreambooth} and LoRA~\cite{hulora2022} serve as baseline methods. All compared techniques employ the loss function prescribed by DreamBooth. For both DreamBooth and LoRA, configurations adhere closely to the respective original papers, utilizing their recommended hyperparameters. Supplemental results can be found in Appendix~\ref{app:settings}, \ref{app:coft_oft}, \ref{app:more_qualitative}, \ref{app:failure}.

\textbf{Finetuning stability and convergence}. We begin by assessing both the stability during finetuning and the rate of convergence for DreamBooth, LoRA, OFT, and COFT. These results are presented in Figure~\ref{fig:teaser} and Figure~\ref{fig:db_convergence}. It is evident that the diffusion model remains stably finetuned when employing either COFT or OFT. After 400 iterations, DreamBooth and both variants of OFT exhibit strong controllability, whereas LoRA fails to maintain subject identity. Upon reaching 2000 iterations, DreamBooth starts producing degenerate (collapsed) images, and LoRA misattributes yellow coloration, generating yellow fur rather than the desired yellow shirt. In contrast, OFT and COFT continue to demonstrate stable and reliable image control throughout. These findings reinforce that our OFT framework attains both rapid and robust convergence in the context of subject-driven generation. An important aspect is the resilience to variation in the number of finetuning iterations, which reduces the need for meticulous adjustment of iteration counts across subjects. With OFT and COFT, one can adopt a comparatively high iteration setting without intensive hyperparameter optimization. Notably, for COFT with an appropriate $\epsilon$, selecting both the learning rate and iteration number becomes straightforward.

\textbf{Quantitative comparison}. Using the protocol established in \cite{ruiz2023dreambooth}, we quantitatively evaluate subject fidelity using DINO~\cite{caron2021emerging} and CLIP-I~\cite{radford2021learning}, prompt fidelity using CLIP-T~\cite{radford2021learning}, and sample diversity using LPIPS~\cite{zhang2018unreasonable}. CLIP-I is determined as the average pairwise cosine similarity between CLIP embeddings of generated and real images, while DINO employs ViT S/16 DINO embeddings analogously. CLIP-T is defined as the mean cosine similarity between embeddings of the text prompt and generated outputs. To measure diversity, we calculate the average LPIPS cosine similarity among images generated with the same prompt and subject. Results, displayed in Table~\ref{table:dreambooth_final}, indicate that both COFT and OFT surpass DreamBooth and LoRA by significant margins on the DINO and CLIP-I metrics, while delivering comparable or marginally superior results in prompt fidelity and diversity assessments. For rigor, we execute each finetuning method 30 times per pretrained model, each with a distinct random seed, thus reducing the effect of stochasticity. Collectively, these outcomes demonstrate that the OFT framework achieves not only superior stability and convergence, but also delivers more consistent and higher quality final results.

\textbf{Qualitative comparison}. To provide a clearer illustration of the advantages offered by OFT, we present a set of randomly chosen instances showcasing subject-driven generation in Figure~\ref{fig:dreambooth_final}. For equitable evaluation, every example is produced using the same finetuned model for each method, and no specific text prompt is separately optimized to enhance particular results. The optimal model for each approach is selected based on the highest validation CLIP metrics. Reviewing Figure~\ref{fig:dreambooth_final}, it is evident that both OFT and COFT exhibit strong abilities in maintaining the semantic features of the subject. In contrast, LoRA frequently does not preserve subject identity (\emph{e.g.}, the LoRA variant entirely fails to maintain the subject's identity in the bowl case). Additionally, both OFT and COFT demonstrate superior control in adhering to text prompts, whereas DreamBooth, though able to retain the subject, often does not faithfully follow the textual instruction (\emph{e.g.}, the first row in the bowl example). These qualitative findings indicate that our OFT framework outperforms others in simultaneously achieving precise controllability and robust subject preservation. Furthermore, the OFT approach is robust to the number of iterations, achieving strong performance even with many iterations, in contrast to DreamBooth and LoRA, which do not perform well under similar conditions. Additional qualitative illustrations can be found in Appendix~\ref{app:more_qualitative}. Also, a human evaluation, detailed in Appendix~\ref{app:human}, provides further support for the advantages of OFT.

\subsection{Controllable Generation}

\textbf{Settings}. As comparative baselines, we utilize ControlNet~\cite{zhang2023adding}, T2I-Adapter~\cite{mou2023t2i}, and LoRA~\cite{hulora2022}. Three difficult controllable generation benchmarks are addressed in this work: Canny edge to image~(C2I) on the COCO dataset~\cite{lin2014microsoft}, segmentation map to image~(S2I) on the ADE20K dataset~\cite{zhou2017scene}, and landmark to face~(L2F) on the CelebA-HQ dataset~\cite{karras2018progressive,xia2021tedigan}. All methods are used to finetune Stable Diffusion~(SD) v1.5 across the three datasets, each for 20 epochs. Supplementary results appear in Appendix~\ref{app:more_qualitative},\ref{app:more_control_task},\ref{app:failure}.

\textbf{Convergence}. We analyze the rate of convergence for ControlNet, T2I-Adapter, LoRA, and COFT on the L2F benchmark, using both quantitative and qualitative approaches. For our quantitative assessment, the principal metric is the mean $\ell_2$ distance between the control face landmarks and those generated by the model. As illustrated in Figure~\ref{fig:face_convergence}, we chart the evolution of face landmark error as each method is finetuned across varying numbers of epochs. Our findings indicate that COFT and OFT converge considerably faster compared to their counterparts. Notably, LoRA requires 20 epochs to attain the same performance level as OFT achieves by its 8th epoch. Importantly, both OFT and COFT have a comparable count of trainable parameters to LoRA (and far fewer than ControlNet), yet they reach convergence more efficiently than the other approaches examined. Additionally, the efficiency of OFT’s convergence is further corroborated by results in Figure~\ref{fig:teaser}: the example on the right demonstrates OFT’s greater data efficiency relative to ControlNet and LoRA, being able to converge robustly with only 5\% of the ADE20K dataset. For our qualitative analysis, we compare OFT, COFT, and ControlNet, as ControlNet most closely matches our landmark error. Visualizations in Figure~\ref{fig:face_convergence_qual} confirm that both OFT and COFT not only converge reliably, but their generated face poses increasingly match the control landmarks over training. Unlike the steady, gradual improvement seen with OFT and COFT, ControlNet’s controllability appears abruptly after the 8th epoch, a pattern consistent with the abrupt convergence observed in \cite{zhang2023adding}. This consistent convergence in OFT leads to more stable and foreseeable training trajectories, facilitating practical use and simplifying the search for optimal hyperparameters.

\textbf{Quantitative comparison}. To assess the effectiveness of controllable generation, we propose a control consistency metric. This metric operates by extracting the control signal from the generated output and measuring its agreement with the original input control. In the C2I scenario, performance is assessed using IoU and F1 score, while in the S2I setting, we report mean IoU alongside both mean and overall accuracy. For the L2F evaluation, the principal metric is the mean $\ell_2$ distance computed between the predicted and control landmarks. Comprehensive descriptions of these metrics are provided in Appendix~\ref{app:settings}. All models are evaluated under their respective optimal hyperparameter configurations. According to the results in Table~\ref{table:control_gen}, both OFT and COFT exhibit notably superior and more precise controllability compared to baseline methods. It is also evident that adapter-style techniques (\emph{e.g.}, T2I-Adapter and ControlNet) demonstrate slower convergence and achieve suboptimal results overall. LoRA outperforms ControlNet on the S2I task but underperforms in the C2I and L2F cases. We generally observe that LoRA’s maximum achievable performance remains limited despite rigorous hyperparameter optimization. In contrast, OFT demonstrates continued gains in performance as the parameter count increases, with no apparent saturation point thus far. Importantly, our approach to quantitative evaluation of controllable generation is itself a contribution, as it facilitates accurate measurement of the control abilities of finetuned models, constrained only by the precision of pre-existing segmentation or detection tools.

\textbf{Qualitative comparison}. In addition to quantitative analysis, we present a qualitative evaluation of OFT and COFT against leading approaches such as ControlNet, T2I-Adapter, and LoRA. As illustrated by the randomly generated results in Figure~\ref{fig:control_final}, OFT and COFT demonstrate not only superior realism and fidelity in image synthesis, but also enable precise controllability. For the S2I task, it is evident that LoRA is unable to adhere to the input segmentation map, while ControlNet, OFT, and COFT achieve successful guided generation. Compared to ControlNet, both OFT and COFT consistently deliver more visually detailed and structurally coherent outputs, all with significantly reduced parameter counts. For the C2I task, OFT and COFT are capable of constructing semantically relevant images from basic Canny edge inputs, a feat where T2I-Adapter and LoRA exhibit markedly poorer performance. In the L2F task, our method ensures the highest accuracy in pose-guided face synthesis, even with challenging orientations. Across all three control scenarios, our findings indicate that OFT and COFT surpass state-of-the-art alternatives in image quality, underscoring the strong controllability of the OFT architecture. To further support this qualitative analysis, additional comparative visualizations for these control tasks are included in Appendix~\ref{app:control_exp}, and we extend the demonstration of OFT to other control scenarios (such as dense pose-to-human body, sketch-to-image, and depth-to-image) in Appendix~\ref{app:more_control_task}.

\section{Concluding Remarks and Open Problems}

Inspired by the key role of angular relationships among neurons in shaping visual semantics, we introduce orthogonal finetuning, a straightforward yet powerful method for refining text-to-image diffusion models. Our approach specifically addresses two use cases: subject-driven generation and controllable generation. In comparison with previous approaches, OFT achieves more precise control and greater stability during finetuning, all while utilizing fewer tunable parameters. Crucially, OFT maintains computational efficiency during inference, facilitating practical deployment without added computational cost.

The development of OFT gives rise to several compelling questions for future work. To begin with, OFT maintains orthogonality through Cayley parametrization, which entails computing a matrix inverse—a procedure that poses some challenges for scaling the method. Although we employ block diagonal parametrization as a workaround, finding ways to accelerate the differentiable computation of the matrix inverse is still an open issue. Another avenue for exploration stems from the compositional properties of OFT: since orthogonal matrices from different OFT finetuning tasks can be combined via multiplication and still yield an orthogonal matrix, it is intriguing to investigate whether this operation preserves the learned information from all individual downstream tasks. Lastly, the parameter efficiency of OFT hinges on the block diagonal structure, which can introduce additional biases and limit the adaptability of the model. Discovering more effective strategies to enhance parameter efficiency without these drawbacks remains a pressing open challenge.

\end{document}